\section{A hierarchy of geometric correlation scales}
\label{sec:hierarchy}

The intervention data allow a scale-resolved statement stronger than a
single endpoint comparison.
For each \(g\), we subtract the exact finite-\(D\) Jacobi curve and test the
entire residual process with a simultaneous band.
For the physical ensemble we repeat the same procedure using eight
independently seeded blocks.
Long-range rigidity is tested independently through the number variance
\begin{equation}
 \Sigma^2(L)=
 \left\langle\bigl[N_\xi(x,L)-L\bigr]^2\right\rangle ,
 \label{eq:number-variance}
\end{equation}
where \(N_\xi(x,L)\) counts unfolded curvature levels in a sliding interval of
length \(L\).

\begin{figure*}[t]
 \includegraphics[width=\textwidth]{figure_4_geometric_hierarchy_v2.pdf}
 \caption{
 \textbf{Controlled hierarchy from structured to random projector geometry.}
 (a) Residual connected SFF
 \(K_{F,c}(\tau;g)-K_{J,c}^{(50)}(\tau)\) along the tangent-scrambling
 interpolation.
 (b) Root-mean-square residual over the preregistered range
 \(0.25\leq\tau\leq1.5\).
 The local adjacent-gap statistic reaches the Haar interval at
 \(g=\GeometricLocalOnset\), while the first nontrivial registered SFF
 compatibility window occurs at \(g=\GeometricRampOnset\).
 (c) Physical random-local-tangent residual with a simultaneous 95\% band;
 it is compatible with the exact finite-\(D\) curve throughout the registered
 range.
 (d) Physical-minus-Haar number variance.
 Short windows are compatible through \(L=\NumberVarianceExtent\), but the
 positive long-range residual is resolved thereafter.
 Onset lines summarize confidence sets on the sampled grids.
 }
 \label{fig:geometric-hierarchy}
\end{figure*}

Three conclusions follow from Fig.~\ref{fig:geometric-hierarchy}.
First, a small interpolation away from Fourier structure generates local
level repulsion, while the two-level correlation curve becomes active at a
larger interpolation strength.
The registered local and ramp onsets are
\begin{equation}
 g_{\rm local}=\GeometricLocalOnset,\qquad
 g_{\rm ramp}=\GeometricRampOnset .
 \label{eq:g-onsets}
\end{equation}
Second, at the physical random-local endpoint the residual band contains zero
for every registered \(\tau\in[0.25,1.5]\).
This is a curve-level finite-\(D\) test; it is more restrictive than matching
the value or slope at one point.

Third, long-range rigidity defines an additional scale.  The physical number
variance follows the Haar reference through \(L=\NumberVarianceExtent\), and
at \(L=8\) Eq.~\eqref{eq:number-residual} resolves a positive simultaneous
interval.
The curvature point process therefore loses information in the ordered
sequence
\begin{equation}
 \begin{aligned}
 \text{local }P(r)
 &\ \prec\ \text{finite-\(D\) SFF ramp}\\
 &\ \prec\ \text{long-range }\Sigma^2(L)
 \end{aligned}
 \label{eq:geometric-correlation-hierarchy}
\end{equation}
where \(A\prec B\) means that \(A\) becomes compatible with the invariant
reference at a weaker scrambling or shorter correlation scale.

Equation~\eqref{eq:geometric-correlation-hierarchy} defines a geometric
Thouless hierarchy in the precise sense of confidence-resolved correlation
lengths in an unfolded curvature spectrum.  Mapping these geometric scales to
transport, scrambling, and real-time diffusion is a promising dynamical
extension.

Our geometric-ETH statement is operational and complementary to
energy-resolved ETH: after removal of a deterministic
metric envelope, local and intermediate-range correlations of the
projector-response spectrum approach those of a finite invariant point
process.
The held-out one-point covariance improvement and the surviving number
variance show that higher Grassmannian information remains measurable.  The
result is covariance-deformed geometric ergodicity.
